\section{What happened when families saw specific schools?}\label{sec:field}

In Chile's 2021 admissions cycle, \citet{schoolSearch} randomized what families saw on a school map. One group saw the basic map. A second also learned how many nearby schools fell into different price and quality bands. A third saw the same counts plus green pins identifying inexpensive schools with medium or high official quality ratings, alongside their prices and ratings. Families remained free to apply anywhere. This design lets us ask whether named recommendations changed applications, whether more pupils got a seat, and where they enrolled.

\begin{table}[ht]\centering\small
\caption{Changes in the School Explorer trial, in percentage points. Both columns compare families assigned the highlighted map with a randomized comparison group. For the basic map, parentheses contain published standard errors. For the counts-only map, brackets give conservative approximate 95\% intervals calculated from the two published standard errors (Appendix~\ref{sec:field-audit}). Sample sizes vary by outcome. Source: \citet{schoolSearch}, Tables I and A.XVIII.}\label{tab:field}
\begin{tabular}{@{}lrrr@{}}\toprule
Outcome & Basic map & Counts-only map & Pupils\\\midrule
Submitted an application & $+3.0\ (1.3)$ & $+2.3\ [-2.8,7.4]$ & 3,111\\
First-ranked school qualified & $+2.5\ (2.1)$ & $+2.4\ [-5.8,10.6]$ & 2,770\\
Second-ranked school qualified & $+5.4\ (2.2)$ & $+3.8\ [-4.8,12.4]$ & 2,580\\
Received any assignment & $+0.2\ (1.1)$ & $-1.0\ [-5.3,3.3]$ & 2,817\\
Enrolled at assigned school & $+3.5\ (1.9)$ & $+2.0\ [-5.6,9.6]$ & 2,748\\
Enrolled at a qualifying school & $+3.3\ (2.3)$ & $+2.0\ [-7.0,11.0]$ & 2,446\\\bottomrule
\end{tabular}
\end{table}

Among the two ranking outcomes, the clearest response was at the second rank. The first-rank estimate was smaller and its interval included zero. The assignment estimate allows both a loss and a gain (95\% interval $[-2.0,2.4]$ points). Holding the price and quality counts fixed gives the sign pattern of more qualifying second choices and fewer assignments. Both intervals include zero. The rank and enrollment outcomes tell us whether a school met the highlighting criteria. They do not identify whether that particular school appeared on that family's map. The enrollment response also varied across families. Relative to the basic map, enrollment at a qualifying school rose by an estimated 14.6 points for pupils whose mothers completed college and 0.0 points for others (reported equality test $p=.009$). An evaluator looking only at the pooled 3.3 points would miss this difference.

\paragraph{Clearing the same year and grades}
We reran both frozen predictors for the full 2021 applicant pool and counted outcomes for the entry grades targeted in the trial. We matched all 461,223 historical course or no-offer records. Among 200,591 entry-grade pupils linked to the April register, the rank-informed menus covered 160,820 eventual schools versus 160,117 under the rank-blind menus. They placed 183,325 versus 183,634 pupils. Only 6,133 pupils in this group received a different list, yet 237 of the 309 net lost placements fell on pupils whose own list stayed identical. An evaluation restricted to families whose display changed would therefore miss most of the loss. Appendix~\ref{sec:field-audit} estimates own-menu effects in 640 national allocations and gives a three-pupil example in which an individual test and a full rollout have opposite signs. This comparison holds submitted choices fixed across policies. Public records identify the national cohort, but contain neither trial arms nor the school IDs displayed to families. The placement cost of deploying either predictor remains a counterfactual.

A separate Chilean trial measured whether families actually enrolled at a school on a card naming nearby options. Among 975 families yet to enroll, assignment to the card-and-video intervention reduced enrollment at a listed school by 15.9 points (SE 3.3). The schools these families chose had higher estimated value added \citep{informedEquilibrium}. This is a reason to record both what a recommender predicted before deployment and where families chose to enroll after seeing it.
